\documentclass{article}
\usepackage[margin=2cm]{geometry}
\usepackage{amsmath}
\usepackage{graphicx}
\usepackage{booktabs}
\usepackage[utf8]{inputenc}
\begin{document}
The expectation values of the operator \(\mathcal{O}_{\phi^{3}}\) and the trace of the energy-momentum tensor are independent of \(\Omega\). As a check, we note that the four one-point functions satisfy the following operator relation, which is associated to the violation of conformal invariance by non-zero classical beta functions,
\[
\left\langle T_{i}^{i}\right\rangle =-\sum_{\mathcal{O}}\left(d-\Delta_{\mathcal{O}}\right)\phi_{\mathcal{O}}\left\langle\mathcal{O}\right\rangle
\]
\subsection*{0.0.1 Derivative of the free energy}
Following , we may now compute the derivative of \(F\) with respect to \(\alpha\) as follows. First we note that
\[
\left.\frac{d F}{d \alpha}=\frac{d S_{\text {ren }}}{d \alpha}=\lim _{\epsilon \rightarrow 0} \int d^{5} x \sum_{\text {fields } \Phi} \frac{\delta\left(\sqrt{\gamma} \mathcal{L}_{\text {ren }}\right)}{\delta \Phi} \left.\frac{d \Phi}{d \alpha}\right|_{z=\epsilon}
\]
In our case, the terms appearing in the sum over fields are
\[
\begin{array}{l}
\frac{\delta\left(\sqrt{\gamma} \mathcal{L}_{\text { ren }}\right)}{\delta \sigma}=\sqrt{\gamma}\left\langle O_{\sigma}\right\rangle \epsilon^{3}+\ldots \quad \frac{\delta\left(\sqrt{\gamma} \mathcal{L}_{\text{ren }}\right)}{\delta \delta \phi^{0}}=\sqrt{\gamma}\left\langle O_{\theta}^{0}\right\rangle \epsilon^{4}+\ldots \\
\frac{\delta\left(\sqrt{\gamma} \mathcal{L}_{\text {\ren }}}\right)}{\delta \phi^{3}}=\sqrt{\gamma}\left\langle O_{\phi}^{3}\right\rangle \epsilon^{3}+\ldots \quad \frac{\delta\big(\sqrt{\gamma} \mathcal{L}_{\text {ren }}}{\delta \gamma^{i j}}=\frac{1}{2} \sqrt{\gamma}\left\langle T_{i j}\right\rangle \epsilon^{5}+\ldots
\end{array}
\]
The dots represent terms of strictly lower order in \(\epsilon\). Furthermore, from the form of the UV asymptotic expansions, we have
\[
\begin{array}{l}
\frac{d \sigma}{d \alpha}=\frac{3}{4} \alpha \epsilon^{2}+O\left(\epsilon^{3}\right) \quad \frac{d \phi^{0}}{d \alpha}=\epsilon+O\left(\epsilon^{3}\right) \\
\frac{d \phi^{3}}{d \alpha}=\left(1-\alpha \frac{d f_{k}}{d \alpha}\right) e^{-f_{k}} \epsilon^{2}+O\left(\epsilon^{3}\right) \quad \quad \frac{d \gamma^{i j}}{d \alpha}=-2 \frac{d f_{k}}{d \alpha} e^{-2 f_{k}} \epsilon^{2}+O\left(\epsilon^{2}\right)
\end{array}
\]
Combining the pieces, with the results for the one-point functions in, we find that the contribution of the metric in is suppressed by \(\epsilon^{2}\) compared to other terms. The derivative of the free energy is then
\[
\begin{array}{l}
\frac{d F}{d \alpha}=\lim _{\epsilon \rightarrow 0} \int d^{5}
x \sqrt{\gamma} \, \epsilon^{5} \, \left[ \frac{3}{2} \beta
\mathrm{e}^{-f_{k}} + \frac{1}{2} \beta
\mathrm{e}^{-f_{k}} \, \left(1-\alpha \frac{d f_{k}}{d \alpha}\right)
+ O(\epsilon) \right] \\
\quad \quad \quad \quad \quad \quad \quad \quad \quad \quad \quad
\quad \quad \quad \quad \quad \quad \quad \quad \quad \; \quad \quad \quad \quad \quad \quad \quad \quad \quad \; = \mathrm{vol}_{0} \, \left( S^{5}\right) \, \frac{1}{2} \beta
\, e^{4 f_{k}} \, \left(4-\alpha \frac{d f_{k}}{d \alpha}\right)
\end{array}
\]
where \(\operatorname{vol}_{0}\left(S^{5}\right)=\pi^{3}\) is the volume of a unit \(S^{5}\). The \(\Omega\) dependence in the one-point functions cancels out, consistent with the fact that \(F\) itself is independent of \(\Omega\). We thus obtain the final result
\[
\frac{d F}{d \alpha}=\frac{\pi^{2}}{8 G_{6}} \beta \, e^{4 f_{k}} \, \left(4-\alpha \frac{df_{k}}{d \alpha}\right)
\]
\end{document}
